% =========================
% PHỤ LỤC
% =========================

\chapter*{PHỤ LỤC}
\addcontentsline{toc}{chapter}{PHỤ LỤC}
\markboth{PHỤ LỤC}{PHỤ LỤC}

% -------------------------------------------------
% Bảng phân chia nhiệm vụ
% -------------------------------------------------
\section*{Bảng phân chia nhiệm vụ}
\phantomsection
\addcontentsline{toc}{section}{Bảng phân chia nhiệm vụ}

\begin{table}[H]
\centering
\caption{Bảng phân chia nhiệm vụ các thành viên}
\label{tab:task-assignment}

\renewcommand{\arraystretch}{1.4}
\begin{tabular}{
|>{\centering\arraybackslash}p{4cm}
|>{\centering\arraybackslash}p{8cm}
|>{\centering\arraybackslash}p{3cm}|
}
\hline
\textbf{Tên thành viên} & \textbf{Nhiệm vụ} & \textbf{Tỉ lệ hoàn thành} \\ \hline

Nguyễn Tiến Phát &
\begin{itemize}
    \item Thiết kế cơ sở hạ tầng hệ thống
    \item Viết script Terraform
    \item Cấu hình Ansible cho các thành phần hạ tầng
    \item Triển khai và cấu hình Kubernetes
    \item Chuẩn bị slide thuyết trình
    \item Viết báo cáo
\end{itemize}
& 100\% \\ \hline

Nguyễn Tài Phú &
\begin{itemize}
    \item Triển khai ứng dụng microservices
    \item Viết và cấu hình Kubernetes manifest cho ứng dụng
    \item Viết báo cáo
\end{itemize}
& 100\% \\ \hline

\end{tabular}
\end{table}

\clearpage

% -------------------------------------------------
% Tài nguyên đồ án
% -------------------------------------------------
\section*{Tài nguyên đồ án}
\phantomsection
\addcontentsline{toc}{section}{Tài nguyên đồ án}

Nhằm phục vụ cho việc tham khảo, đánh giá và tái sử dụng kết quả của đề tài, nhóm
cung cấp các tài nguyên liên quan bao gồm slide thuyết trình và mã nguồn triển khai
hệ thống. Các tài nguyên này được lưu trữ trực tuyến và có thể truy cập thông qua
các liên kết sau:

\begin{itemize}
    \item \textbf{Slide thuyết trình:} \\
    \href{https://drive.google.com/drive/folders/1DX2yMvPsaUAYSylsWpX3x7XvMinTPM6z}
    {Google Drive – Slide thuyết trình đồ án}

    \item \textbf{Mã nguồn đồ án:} \\
    \href{https://github.com/NT533-Q12-Distributed-Computing}
    {GitHub – NT533-Q12-Distributed-Computing}
\end{itemize}

Các tài nguyên trên phản ánh đầy đủ quá trình thiết kế, triển khai và đánh giá hệ
thống trong phạm vi đề tài, đồng thời có thể được sử dụng làm cơ sở tham khảo cho
các nghiên cứu và đồ án liên quan trong tương lai.

\clearpage

% -------------------------------------------------
% Tài liệu tham khảo
% -------------------------------------------------
\section*{Tài liệu tham khảo}
\phantomsection
\addcontentsline{toc}{section}{Tài liệu tham khảo}

Trong quá trình thực hiện đề tài, nhóm đã tham khảo các tài liệu chính thống
liên quan đến DevOps, hệ thống phân tán, Kubernetes, hạ tầng cloud-native và
các công cụ quan sát hệ thống. Các tài liệu này đóng vai trò nền tảng lý thuyết
và hướng dẫn triển khai cho toàn bộ nội dung của đồ án.

\begin{itemize}
    \item \textbf{Kubernetes và hệ thống phân tán}
    \begin{itemize}
        \item \href{https://kubernetes.io/docs/concepts/}
        {Kubernetes Documentation – Core Concepts}
        
        \item \href{https://kubernetes.io/docs/tasks/run-application/horizontal-pod-autoscale/}
        {Kubernetes Documentation – Horizontal Pod Autoscaler}
        
    \end{itemize}

    \item \textbf{Giám sát và quan sát hệ thống (Observability)}
    \begin{itemize}
        \item \href{https://prometheus.io/docs/introduction/overview/}
        {Prometheus – Monitoring System Overview}
        
        \item \href{https://grafana.com/docs/}
        {Grafana – Visualization and Dashboard Documentation}
        
        \item \href{https://grafana.com/oss/loki/}
        {Grafana Loki – Log Aggregation System}
        
        \item \href{https://www.cncf.io/projects/}
        {Cloud Native Computing Foundation – Cloud Native Projects}
    \end{itemize}

    \item \textbf{Hạ tầng và tự động hoá}
    \begin{itemize}
        \item \href{https://developer.hashicorp.com/terraform/intro}
        {HashiCorp Terraform – Infrastructure as Code}
        
        \item \href{https://docs.ansible.com/}
        {Ansible Documentation – Automation Platform}
        
        \item \href{https://docs.aws.amazon.com/eks/latest/userguide/what-is-eks.html}
        {Amazon EKS – User Guide}
    \end{itemize}

    \item \textbf{Độ tin cậy hệ thống}
    \begin{itemize}
        \item \href{https://sre.google/books/}
        {Google – Site Reliability Engineering}
    \end{itemize}
\end{itemize}
